\section{The first two arches: groupoids and the Ore condition}\label{sec:arches}

Whether $\PSh(C)$ is Boolean or De Morgan is a question about all presheaves on $C$, of which there are infinitely many. The two bridges of this section reduce it to a finite check on the shape $C$ itself.

\subsection{Boolean presheaf topoi}

\begin{bridge}[Boolean]\label{br:bool}
$\PSh(C)$ is Boolean iff $C$ is a groupoid.
\end{bridge}

\begin{proof}
Suppose $C$ is a groupoid, $A\subseteq F$, and $x\in F(U)$ not in $A$. If some restriction $\res xf$ were in $A$, so would be $x = \res{(\res xf)}{f^{-1}}$. So $x\in\neg A$, and $A\cup\neg A = F$.

Conversely, suppose $\PSh(C)$ Boolean and let $f : V\to U$. Take $F = \y U$ and $A$ = the sieve generated by $f$ (the arrows $f\circ h$). The element $\id_U$ is in $A$ or in $\neg A$. It is not in $\neg A$, because its restriction along $f$ is $f\in A$. So $\id_U\in A$: $\id_U = f\circ h$ for some $h : U\to V$. Every arrow therefore has a right inverse. Applying this to $h$, there is $k$ with $h\circ k = \id_V$; then $f = f\circ h\circ k = k$, so $h\circ f = \id_V$ and $f$ is invertible.
\end{proof}

\begin{example}
$G$-sets form a Boolean topos. So do presheaves on any equivalence relation, and \q{sets with an involution} (presheaves on $\B(\mathbb Z/2)$). Graphs, automata over a non-empty alphabet, simplicial sets, the topos of trees are not Boolean: their shapes have non-invertible arrows.
\end{example}

\subsection{The Ore condition}

\begin{definition}[Ore condition]
A category $C$ satisfies the \emph{(right) Ore condition} if every cospan can be completed to a commutative square: for all $f : Y\to X$ and $g : Z\to X$ there are $W$, $h : W\to Y$, $k : W\to Z$ with $f\circ h = g\circ k$.
\[
\begin{tikzcd}
W \arrow[r,dashed,"k"] \arrow[d,dashed,"h"'] & Z \arrow[d,"g"]\\
Y \arrow[r,"f"'] & X
\end{tikzcd}
\]
\end{definition}

The square is not required to be a pullback, nor unique: any completion will do.

\begin{example}[When Ore holds]\label{ex:orehold}
\begin{enumerate}
\item A groupoid: take $W = Y$, $h = \id$, $k = g^{-1}\circ f$.
\item A category with an \emph{initial object} $0$: take $W = 0$ (the two composites $0\to X$ agree, since there is only one). So $\Delta_+$, $\mathsf{Fin}$, $\mathsf{FI}$ satisfy Ore.
\item A category with pullbacks: take the pullback.
\item A preorder in which any two elements below a common element have a common lower bound. For instance a chain such as $\omega$ ($W=\min(Y,Z)$), or any lattice.
\item $\B M$ for a monoid $M$: a cospan is a pair $a,b\in M$, a square is a pair $c,d$ with $ac = bd$. So $\B M$ satisfies Ore iff $aM\cap bM\neq\emptyset$ for all $a,b$. This holds if $M$ is commutative ($ab = ba$) or has a zero element ($a0 = b0$).
\end{enumerate}
\end{example}

\begin{example}[When Ore fails]\label{ex:orefail}
\begin{enumerate}
\item The graph shape $\mathsf G$: the cospan $V\xrightarrow{s}E\xleftarrow{t}V$. The only arrows into $V$ are from $V$ itself (identities), and $s\neq t$.
\item $\Delta$: the cospan $[0]\xrightarrow{\delta^1}[1]\xleftarrow{\delta^0}[0]$ picking the two vertices. Any $h : [n]\to[0]$ is constant, so $\delta^1\circ h$ is constantly $0$ and $\delta^0\circ k$ constantly $1$.
\item $\mathsf{Fin}_{>0}$ and the cube category, for the same reason: two distinct points of a set cannot be reached by maps from a non-empty set that agree.
\item The free monoid $\Sigma^*$ on $\Sigma\ni x,y$ with $x\neq y$: every element of $x\Sigma^*$ starts with $x$, every element of $y\Sigma^*$ with $y$.
\item The \q{V} preorder $\{a<b,\ a<c\}$ seen with arrows reversed, i.e.\ the preorder $\{b<a,\ c<a\}$ (the \q{$\Lambda$} shape): $b\to a\leftarrow c$ has no common lower bound.
\end{enumerate}
\end{example}

\subsection{De Morgan presheaf topoi}

\begin{bridge}[De Morgan]\label{br:dm}
$\PSh(C)$ is De Morgan iff $C$ satisfies the Ore condition.
\end{bridge}

\begin{proof}
Suppose $C$ satisfies Ore. By Lemma~\ref{lem:dmpointwise}, we must show: if some restriction $\res xf$ of $x\in F(X)$ lies in $A$ ($f : Y\to X$), then every restriction $\res xg$ ($g : Z\to X$) has a further restriction in $A$. Complete $f,g$ to a square $f\circ h = g\circ k$. Then $\res{(\res xg)}{k} = \res x{g\circ k} = \res x{f\circ h} = \res{(\res xf)}{h}\in A$, since $A$ is closed under restriction.

Conversely, suppose $\PSh(C)$ De Morgan, and let $f : Y\to X$, $g : Z\to X$. Take $F = \y X$, $A$ = the sieve generated by $f$, and the element $\id_X$. Its restriction along $f$ is $f\in A$, so $\id_X\notin\neg A$; by De Morgan, $\id_X\in\neg\neg A$. So its restriction $g$ has a further restriction in $A$: $g\circ k = f\circ h$ for some $h$, $k$. That is the square.
\end{proof}

\begin{remark}
The two proofs have the same structure, and it is worth remembering: \emph{to test a property of all presheaves, test it on the representable presheaves and their sieves}. Representables are the \q{free} presheaves; the sieve generated by an arrow is the smallest sub-presheaf containing it.
\end{remark}

\begin{example}
By Example~\ref{ex:orehold} and Bridge~\ref{br:dm}: augmented simplicial sets, presheaves on $\mathsf{Fin}$ or $\mathsf{FI}$, the topos of trees, and $M$-sets for a commutative monoid or a monoid with zero form De Morgan topoi. By Example~\ref{ex:orefail}: graphs, simplicial sets, symmetric simplicial sets, cubical sets, and automata over two or more letters do not. Example~\ref{ex:Delta1} and Figure~\ref{fig:edge} are exactly the proof of Bridge~\ref{br:dm} run on the failing cospans.
\end{example}

\begin{exercise}\label{exo:oreproduct}
Show that if $C$ and $D$ satisfy Ore, so does $C\times D$. Deduce that bisimplicial sets (presheaves on $\Delta\times\Delta$) are not De Morgan, but that presheaves on $\omega\times\omega$ are.
\end{exercise}

\subsection{The bridges in Lean}

Bridges~\ref{br:bool} and~\ref{br:dm} are classical \cite[D4.6.3]{elephant}; see also \cite[Prop.~14]{caramello-chamoun} and, for monoids, \cite{hemelaer-rogers}. They are proved in Lean in \lp{}:
\leanbox{\lean{isBooleanTopos_presheaf_iff_isGroupoid}, \lean{isDeMorganTopos_presheaf_iff_rightOreCondition}. The Lean proofs follow the same route as above: \lean{isDeMorganTopos_presheaf_iff} is Lemma~\ref{lem:dmpointwise}, and the converse uses the sieve generated by an arrow, as a sub-presheaf of a representable.}
